\section{\texorpdfstring{Определение области. Определение функции
комплексного переменного, заданной на области
\(G \subset \mathbb{C}\).}{Определение области. Определение функции комплексного переменного, заданной на области G \textbackslash subset \textbackslash mathbb\{C\}.}}\label{ux43eux43fux440ux435ux434ux435ux43bux435ux43dux438ux435-ux43eux431ux43bux430ux441ux442ux438.-ux43eux43fux440ux435ux434ux435ux43bux435ux43dux438ux435-ux444ux443ux43dux43aux446ux438ux438-ux43aux43eux43cux43fux43bux435ux43aux441ux43dux43eux433ux43e-ux43fux435ux440ux435ux43cux435ux43dux43dux43eux433ux43e-ux437ux430ux434ux430ux43dux43dux43eux439-ux43dux430-ux43eux431ux43bux430ux441ux442ux438-g-subset-mathbbc.}

\textbf{Определение:} Окрестность точки \(z_{0}\): \[
U_{\delta}(z_{0})=\{ z\in \mathbb{C}:|z-z_{0}|<\delta \}
\] \textbf{Определение:} Окрестность \(\infty\): \[
U_{R}(\infty)=\{ z\in \mathbb{C}:|z|>R \}
\]

\textbf{Определение:} Множество \(G \subset \overline{\mathbb{C}}\)
называется связным, если \(\forall z_1, z_2 \in G\ \exists\) непрерывная
кривая \(\gamma \subset G\), соединяющая \(z_{1}\) и \(z_2\).

\textbf{Определение:} Множество \(G\subset \overline{\mathbb{C}}\)
называется открытым, если
\(\forall z_{0}\in G \ \ \exists U(z_{0})\subset G\)

\textbf{Определение:} \(G\) - область в \(\overline{\mathbb{C}}\) - это
открытое, связное множество.

\textbf{Определение:} Функция комплексного переменного:
\(w=f(z), \ z\in G \overset{f}{\to}w\in f(G)\subset \mathbb{C}\). \[
z=x+iy, \quad w=u(x,y)+i\cdot v(x,y), \quad u=\mathrm{Re}f(z), \ v=\mathrm{Im}f(z)
\]
